% ============================================================
\chapter{テキストのセンチメントと LLM}\label{ch:text}

ニュース、掲示板、決算説明会の文章からビューを作る研究は、BERT のような言語モデルの登場で増え、最近は大規模言語モデル（LLM）に直接リターンを予想させる研究も出てきた。本プロジェクトもその 1 つを再現している。この章では、代表的な研究を、ビュー $Q$ と不確かさ $\Omega$ をどう作っているかに注目して読み、そのうえで、LLM の回答のばらつきを $\Omega$ に使うことの是非を考える。

\section{テキストからビューを作る流れ}

どの研究も、おおむね次の流れである。

\begin{enumerate}
\item 文章を集め、銘柄と日付に結びつける。
\item 文章を数値にする（辞書で数える、BERT で分類する、LLM に聞く）。
\item 数値を予想リターンに直す（回帰、LSTM などの予測モデル、価格のシミュレーション、LLM の回答そのもの）。
\item 予想をビュー $Q$ にし、不確かさ $\Omega$ を決める。
\end{enumerate}

研究によって工夫があるのは 2 と 3 で、4 の $\Omega$ は簡単に済ませているものが多い。ところが、BL の結果を大きく決めるのは 4 である（\chapref{ch:overview}）。

\section{BERT のセンチメントをビューにした研究}

\begin{paperbox}[Colasanto, Grilli, Santoro and Villani（2022）, \textit{BERT's sentiment score for portfolio optimization: a fine-tuned view in Black and Litterman model}, Neural Computing and Applications 34]
\begin{itemize}
\item Financial Times の記事のセンチメントの極性の点数を、BL の新しいビューとして使う。
\item 記事が出たあとの数取引時間の価格の経路をモンテカルロ法でシミュレーションし、得られた価格を利回りに直して、時間ごとの価格に基づく「動的な」最適化のビューにする。
\item センチメントを使わないビュー（初めのポートフォリオを変えない）と比べる。
\end{itemize}
本文は確認できず、要旨の範囲で記した。
\end{paperbox}

\begin{paperbox}[Sun, Fu and Liu（2026）, \textit{Investor sentiment and stock portfolio construction: A study based on the Black-Litterman model}, PLOS ONE 21(9)]
\begin{itemize}
\item \textbf{データ}：上海 50 指数の 8 銘柄（保利発展、中国聯通、恒瑞医薬、伊利、中国建築、中国銀行、中国石化、中微公司）について、投資家の掲示板（東方財富の股吧）の書き込みを 2022 年 7 月〜2024 年 12 月に集め、自動投稿や重複を除いて 387,953 件。
\item \textbf{センチメント}：辞書の方法と、中国語 BERT を 3 区分（悪い・中立・良い）に微調整した方法。閲覧数と返信数で重み付けして、銘柄ごとの日次の指数にする。BERT の正解率は、学習に使った領域の検証データで 74.69\%、掲示板の書き込み 1,000 件を人が付けた正解と比べると 52.3\% だった。
\item \textbf{予想}：始値・高値・安値・終値・出来高とセンチメントを LSTM に入れ、対数リターンを予想する。2022 年 7 月〜2024 年 7 月で学習、2024 年 7 月 11 日〜12 月 31 日で検証。
\item \textbf{BL}：LSTM の予想を絶対ビュー $Q$ にし、ビューの不確かさは $8\times8$ の単位行列と書かれている。事前の共分散は過去 200 取引日。毎日見直し、117 取引日で評価。
\item \textbf{結果}：取引費用なしで、BERT のビューの BL は年率 106.58\%、シャープレシオ 2.64、ソルティノレシオ 4.18。辞書とセンチメントなしの BL より良く、取引費用を入れても最も高い年率リターンを保った。
\item \textbf{限界}（著者による）：書き込みは個人投資家が中心、文章の雑音、汎用の BERT の限界、滑りや比率の上限を入れていない、8 銘柄・1 期間だけ。
\end{itemize}
\end{paperbox}

調査メモは、この論文の年率 106.58\%、シャープレシオ 2.64、ソルティノレシオ 4.18 を「極めて優れた」成績として紹介している。数字は論文どおりだが、次の条件を添えて読む必要がある。

\begin{itemize}
\item 評価期間は 117 取引日（約 5 か月半）で、その終わりの 2024 年秋は中国株が大きく上がった時期である。年率に直した 106\% は、5 か月余りの上昇を 1 年に引き延ばした数字である。
\item 銘柄は 8 つだけで、どの銘柄を選ぶかで結果が大きく変わりうる。
\item ビューの不確かさが単位行列だとすると、日次の対数リターン（$10^{-4}$ 程度の分散）に対して $\Omega$ の対角が 1 なので、ビューはほとんど採用されないはずである。単位が違う（パーセントで書いている、など）のか、別の扱いをしているのか、本文の範囲では確認できなかった。$\Omega$ の大きさは結果を左右するので、再現するならここを最初に確かめる必要がある。
\item センチメントの分類器の正解率は、実際の書き込みでは 52.3\%（3 区分でのでたらめは 33\%）だった。
\end{itemize}

\begin{paperbox}[Mantshimuli and Muteba Mwamba（2025）, \textit{Enhancing portfolio optimization with multi-LLM sentiment aggregation: A Black-Litterman integration approach}, Investment Management and Financial Innovations 22(3)]
\begin{itemize}
\item 金融向けに微調整した 3 つの LLM のセンチメントの点数を、LSTM で 1 つの「メタ LLM」の点数にまとめ、BL の投資家のビューにする。
\item \SP{} の銘柄で検証し、年率 31.22\%（時価総額加重は 24.57\%）、シャープレシオ 3.02、オメガレシオ 2.48、ジェンセンのアルファ 1.95\% と報告している。単一の LLM のポートフォリオより良かった。
\item 要旨の範囲では、検証期間や取引費用は分からなかった。
\end{itemize}
\end{paperbox}

\section{対象論文：LLM の予想と回答のばらつき}

\begin{paperbox}[Lee, Kim, Kim, Kim and Lee（2025）, \textit{LLM-Enhanced Black-Litterman Portfolio Optimization}, CIKM 2025 FinAI ワークショップ（arXiv 2504.14345 第 2 版）]
\begin{itemize}
\item \textbf{データ}：2025 年 3 月 26 日時点の \SP{} の時価総額上位 50 銘柄。2024 年 6 月〜2025 年 6 月の 13 か月を、検証用 3 か月と評価用 10 か月に分ける。2 週間ごとに見直す。
\item \textbf{LLM}：Gemma-7B、Qwen-2-7B、LLaMA-3.1-8B、GPT-4o-mini。いずれも学習データの締め切りが評価期間より前。
\item \textbf{質問}：過去 2 週間の日次リターン、セクターと \SP{} の日次リターン、会社の情報を渡し、次の 2 週間の日次の平均リターンを 1 つの数で答えさせる。
\item \textbf{ビュー}：各銘柄について 100 回聞き、回答の平均を $Q$、回答の分散を $\Omega$ の対角にする。$P=I$（全銘柄に絶対ビュー）。
\item \textbf{$\tau$}：検証期間で、平均($\Omega$) ÷ 平均($\Sigma$) を中心に 0.5〜1.5 倍の 5 点を試し、シャープレシオで選ぶ。最適化は $w\T\Sigma w-\lambda w\T r$（$\lambda=0.1$）を、全額投資・空売りなしで最小化。取引費用は比率の変化の 0.1\%。
\item \textbf{結果}：年率の成長率は Qwen 28.11\%、Llama 27.51\%、Gemma 15.90\%、GPT 7.68\%、等金額 19.07\%、平均分散 6.07\%。年率のシャープレシオは Llama 1.23、Qwen 1.06。
\item \textbf{解釈}：LLM ごとに一貫した「投資スタイル」（Gemma は悲観、Qwen は控えめな楽観、Llama は大きく散らばる強い確信）があり、成績はそのスタイルと相場の向きが合うかで決まる。
\end{itemize}
\end{paperbox}

本プロジェクトは、この論文の第 1 版（30 回、1 か月の予想）を、全銘柄・30 か月に広げて再現した（『総合報告書』）。調査メモが紹介している「100 回の反復」は第 2 版の設定である。

論文は、回答の分散を $\Omega$ にすることで「確信度が自動で較正される」と述べている。調査メモはさらに進んで、「LLM の出力の統計的なばらつきが、市場の政治的・定性的な情報の不確かさを直接表している」「人が恣意的に決めていた確信度を、テキストデータの内在的なエントロピーから客観的に導出できる」と評価している。次の節で、この主張を検討する。

\section{回答のばらつきは何を測っているか}

同じ質問を何度も聞いたときの回答の散らばりは、LLM が次の語を確率的に選ぶ仕組み（温度 temperature によるサンプリング）から生まれる。散らばりの大きさは、次のようなもので決まる。

\begin{itemize}
\item \textbf{温度の設定}：温度を 0 に近づければ、回答はほとんど同じになる。本プロジェクトは温度 1.0 で聞いている。温度を変えれば $\Omega$ が変わり、BL の結果が変わる。
\item \textbf{質問の書き方}：出力の形式（数値だけ、桁数）や、渡す数値の単位で、散らばり方が変わる。対象論文の第 2 版は、リターンを 100 倍（パーセント）にして渡している。
\item \textbf{入力の値の散らばり}：本プロジェクトでは、回答の散らばりは、渡した月の値動きの大きさと順位相関 0.73 で連動していた（『総合報告書』問題 3）。値動きが大きい銘柄ほど、回答もばらつく。
\end{itemize}

どれも「予想がどれだけ当たるか」とは別のものである。LLM の確信の表明が実際の正しさと合わないことは、言語モデルの研究でも繰り返し報告されている。例えば Xiong ほか（2024）は、LLM に自分の確信度を言わせると、多くの場合に自信過剰になることを示した。Kadavath ほか（2022）は、選択式の問題では大きなモデルの確率がおおむね較正されていることを示したが、それは答えが決まっている問題の話で、来月の株価のように誰にも分からないものの予想とは事情が違う。

本プロジェクトのデータでは、回答の分散と、実際の予想の外れ（誤差の 2 乗）の比は、銘柄・月ごとの中央値で \V{llm.ratio.med}（四分位で \V{llm.ratio.q25}〜\V{llm.ratio.q75}）だった（\chapref{ch:project}）。回答の散らばりは実際の外れの約 8 分の 1 で、BL はその分だけ予想を信じすぎる。

\begin{notebox}[回答のばらつきの使い方]
回答のばらつきが、外れやすさの「順番」の手がかりになる可能性は残る（ばらつく銘柄ほど外れやすい、など）。しかし大きさは合っていない。使うなら、ばらつきの順番だけを使い、大きさは当たり具合で較正する。例えば、$\omega_i=k\cdot(\text{回答の分散}_i)$ として、倍率 $k$ を過去の予想の外れから決める（\chapref{ch:accuracy}）。
\end{notebox}

\section{確率の分布からビューを作る}

LLM に数値を直接答えさせる代わりに、「強気・中立・弱気」の確率を答えさせる方法もある。Florencio（2026）は、論理の述語を確率で表す「neural predicate」の考え方で、この確率から BL の $P$、$q$、$\Omega$ を作る理論の枠組みを提案している。

\begin{itemize}
\item \textbf{ビューの値}：強気の確率を $p_b$、弱気を $p_r$、中立を $p_n$ とし、強気・弱気のときのリターンを $r_{\mathrm{bull}}$、$r_{\mathrm{bear}}$ とする。$q_i=p_b\,r_{\mathrm{bull}}+p_r\,r_{\mathrm{bear}}+p_n\,\Pi_i$（中立なら均衡のまま）。
\item \textbf{不確かさ}：確率の分布のエントロピー $H$ を使う。$\omega_i=\omega_{\min}+(\omega_{\max}-\omega_{\min})\,H/\log3$。分布が 1 つの区分に集中していれば（$H=0$）$\omega_{\min}$、3 区分に均等なら（$H=\log3$）$\omega_{\max}$。$\omega_{\min}$、$\omega_{\max}$ は $\tau\Sigma_{ii}$ の倍数にして、事前の不確かさと単位をそろえる。
\item \textbf{複数の述語の組み合わせ}：述語ごとの $(q,\omega)$ を精度（$1/\omega$）で加重平均し、述語どうしの食い違い（$q$ の分散）が大きいほど $\omega$ を大きくする。
\end{itemize}

\begin{table}[htbp]
\centering
\caption{確率の分布からビューを作る例（月次 \%）。$r_{\mathrm{bull}}=2\%$、$r_{\mathrm{bear}}=-2\%$、$\Pi_i=0.6\%$、$\sigma_i=8\%$、$\tau=0.05$、$\omega/\tau\sigma^2$ を 1〜20 の範囲でエントロピーに比例させた（倍率はこの資料が例として置いたもの）。}
\label{tab:stance}
\tabin{stance}
\end{table}

\tabref{tab:stance} のとおり、確率が強気に集中した銘柄 A は、エントロピーが小さく採用率が高い。3 区分に均等な銘柄 D は、ビューの値も均衡に近く、採用率も低い。仕組みとしては分かりやすい。

ただし、ここでもエントロピーは「LLM がどれだけ迷っているか」であって、「当たり具合」ではない。$\omega_{\min}$、$\omega_{\max}$ をどう置くかで採用率は自由に変わるので、結局は過去の当たり具合で決める必要がある。LLM の確率そのものが較正されているかも確かめる必要がある。本プロジェクトの実験 G（リスクプレミアムの有無の確率を聞く）では、確率の読み取りは使えるが、なめらかではなかった（『総合報告書』）。

\section{複数のモデルの予想をまとめる}

LLM を複数使えば、同じ銘柄について複数のビューが得られる。Spears、Zohren、Roberts（2023）は、同じ資産への複数のビュー（とその不確かさ）をまとめる「融合」の方法を、BL の文脈で整理した。

\begin{description}[leftmargin=2em,style=nextline]
\item[精度加重（PW）] 外れが独立なら、分散の逆数で重み付けして平均するのが最適。まとめた分散は $(\sum_s\hat\Sigma_s^{-1})^{-1}$ で、情報源の数だけ小さくなる。
\item[共分散の交差（CI）] 外れの相関が分からないとき、逆行列の凸結合 $\hat\Sigma^{-1}=\sum_s\omega_s\hat\Sigma_s^{-1}$（$\sum\omega_s=1$）を使う。相関がどうであっても、不確かさを小さく見積もりすぎない（一貫性がある）が、慎重すぎることがある。
\item[逆の共分散の交差（ICI）] CI より慎重さを抑えつつ一貫性を保つ方法。2 つずつまとめる。
\item[共分散の和集合（CU）] どれかの情報源が誤っている（不確かさを小さく言っている）可能性があるときに使う。まとめた分散が、各情報源の分散に平均のずれを足したものを含むようにする。
\end{description}

彼らは、ARIMA、ブースティング、ガウス過程の 3 つのモデルのビューを、BL--APT（\chapref{ch:shape}）で 1992〜2022 年の米国株に当てはめた。年ごとのシャープレシオの中央値は、融合の方法で PW 0.11、CU 0.16、CI 0.38、ICI 0.48 と並び、相関を考える方法のほうが良かった（ただし \SP{} を持ち続けるほうが全体として優れていた）。

\subsection*{外れの相関が大事}

対象論文の表 5 には、4 つの LLM の予想の誤差の 2 乗の平均（MSE）が載っている。これで精度加重の重みを計算すると\tabref{tab:fusion-w}のようになる。

\begin{table}[htbp]
\centering
\caption{対象論文の表 5 の MSE と、精度加重（$1/\mathrm{MSE}$ に比例）の重み。}
\label{tab:fusion-w}
\tabin{fusion_w}
\end{table}

外れが独立なら、まとめた MSE は \V{fu.pw.indep} になり、いちばん良い Qwen（\V{fu.best}）の 3 分の 1 以下になる。しかし、4 つの予想の外れは独立ではない。同じ過去のデータを見て予想し、同じ実現値と比べているので、実現値の側の予想できない動きは、4 つの外れに共通に入る。

\begin{table}[htbp]
\centering
\caption{外れの相関 $\rho$（すべての組で同じとした）を変えたときの、まとめた MSE。「最適な組み合わせ」は、相関を知っていて、重みの和が 1 の範囲で MSE を最小にしたもの（負の重みも許す）。}
\label{tab:fusion-rho}
\tabin{fusion_rho}
\end{table}

\tabref{tab:fusion-rho} のとおり、相関が 0.8 を超えると、精度加重でまとめた MSE は、いちばん良い単独の予想より悪くなる。相関を知っていれば改善の余地はあるが、そのときの重みは大きな正と負の組み合わせになり、相関の推定の誤差に敏感になる。複数の LLM をまとめるときは、独立を仮定した精度加重を避け、CI のような慎重な方法を使うか、それぞれの予想の過去の当たり具合で重みを決める（相関も過去の外れから測る）のがよい。

\begin{tipbox}[コツ：テキストや LLM からビューを作るとき]
\begin{enumerate}
\item 論文の成績は、評価期間の長さ、銘柄の数、取引費用、$\Omega$ の大きさを確かめてから読む。
\item LLM の回答のばらつきやエントロピーは、外れやすさの順番の手がかりにはなりうるが、大きさは当たり具合で較正する。
\item LLM の「スタイル」（全体の楽観・悲観）は相対ビューで取り除き、銘柄間の差だけを使う（\chapref{ch:shape}）。
\item 複数の LLM をまとめるときは、外れの相関を考える。独立を仮定した精度加重は確信を過大にする。
\item 温度や質問の書き方を変えたら、ばらつきの大きさも変わる。設定を変えるたびに較正し直す。
\end{enumerate}
\end{tipbox}

\begin{keybox}[この章のまとめ]
\begin{itemize}
\item テキストからビューを作る研究は、センチメントや予想の作り方に工夫があり、$\Omega$ は簡単に済ませているものが多い。
\item Sun ほか（2026）の年率 106.58\% は、8 銘柄・117 取引日・$\Omega$ が単位行列という条件の結果である。
\item LLM の回答のばらつきは、温度や質問の書き方で決まり、当たり具合を表さない。本プロジェクトでは実際の外れの約 8 分の 1 だった。
\item 複数の LLM の予想は、外れが強く相関するので、独立を仮定してまとめると確信が過大になる。
\end{itemize}
\end{keybox}
